


\section[Modelos de replicación]{Modelos de replicación del ADN}



\section{Proceso de replicación} % En procariotas

Es el proceso en que el ADN de una célula se replica para preparase para la división celular. Este proceso se compone de 3 partes. En la primera, la de iniciación, proteínas especializadas se unen al sitio de replicación. Luego, en la etapa de elongación, la ADN polimerasa comienza a sintetizar una nueva cadena de ADN. Por último, en la etapa de terminación, la endo/exonucleasas terminan el proceso.

\subsection{Componentes de la replicación}

\begin{itemize}
    \item Topoisomerasa: Esta va  por fuera del bucle. Desenrolla la cadena de ADN para que, posteriormente, la helicasa pueda romperla.
    \item Helicasa: Esta va por dentro del bucle. Rompe los enlaces de hidrógeno que unen a la cadena de ADN.
    \item Fragmentos de Okasaki: Se componen de ARN.
    \item Lugar de unión: Lugar donde las proteínas se unen a la hebra rezagada.
    \item Nucleotidos: Se transforman de ADN a ARN por acción de la proteína ADN ligasa. 
        \begin{itemize}
            \item Procariotas: 1000-2000 nucleótidos.
            \item Eucariotas: Menor cantidad de nucleótidos.
        \end{itemize}
        
    \item Cadenas
        \begin{itemize}
            \item Cadena principal: Trabaja de 3' a 5'
            \item Cadena rezagada: Trabaja de 5' a 3'. Aquí se aplican los fragmentos de Okazaki
        \end{itemize}
    \item Proteínas de unión: Se unen a los puntos de iniciación
    \item Proteínas de unión monocatenarias: Bloquean que se vuelva a enrollar la cadena
    \item ADNpolimerasa (ADNpol): Trabaja de 5' a 3'. Necesita de un grupo hidroxilo libre ($OH^{-}$).
    \item Primers, cebadores o iniciadores: Se colocan por actuar de la ARN primasa. Proporcionan el grupo hidroxilo libre ($OH^{-}$) que la ADNpol necesita. Se componen de 5-10 nucleotidos.
\end{itemize}


\subsection{Iniciación} % Origen de replicación (Ori)
\subsection{Elongación}
\subsection{Terminación}


\section{Particularidades de la replicación eucariota}

Hay una serie de factores que influyen en que la replicación del ADN de células eucariotas sea algo más compleja. Estos son que el ADN de células eucariotas es más largo y se encuentra ``enrredado'' en proteínas llamadas histonas, las cuales al agregarse forman nucleosomas. Al darse la replicación del ADN, la hebra vieja se queda con los nucleosomas viejos y la hebra nueva se ``enrreda'' en histonas nuevas.